% ThemeForge demo slides deck.
% 
% 
% 
% 
\documentclass[11pt, aspectratio=169]{beamer}
\providecommand\ThemeName{}
\providecommand\ThemeOptions{}

% Lines required for standalone compiling within the project files
% \makeatletter
% \def\input@path{{../themes/}}
% \makeatother

\edef\UseTheme{\noexpand\usetheme[\ThemeOptions]{\ThemeName}}\UseTheme

\usepackage{booktabs}
\usepackage{appendixnumberbeamer}
\hypersetup{pdflang=en-GB, pdfdisplaydoctitle=true}

\title[New Beamer]{Modern Slides with Beamer}
\subtitle{A tour of the \ThemeName{} theme}
\author[Rose C.]{Rosario Corazon}
\institute{ThemeForge}
\date{\today}

\begin{document}

% Title frame. 
% `plain` to keep normal frame elements interfering with the title.
% `noframenumbering` so the frame counter skips the title page
\begin{frame}[plain, noframenumbering]
    \titlepage    
\end{frame}


\begin{frame}{Agenda}
    \tableofcontents
\end{frame}

\section{Typography}

\begin{frame}{Text and Emphasis}{Readable from the back of the room}
    A good slide carries one idea. Body text is set large, with a comfortable line length and \alert{alerted text} for the one takeaway. \emph{Emphasis} stays subtle, and \texttt{monospace} for code.

    \medskip
    Inline maths sits on the baseline: $E_n = (n + \frac{1}{2})\hbar\omega$, and displays breathe:
    \[
        i\hbar \frac{\partial}{\partial t}\Psi(\mathbf{r}, t) = \left[ -\frac{\hbar^2}{2m}\nabla^2 + V(\mathbf{r}, t) \right] \Psi(\mathbf{r}, t)
    \]
\end{frame}

\begin{frame}{Lists}
    
\end{frame}


\begin{frame}{Blocks}
    \begin{block}{Definition}
        A \emph{theme} controls the colours, fonts and layout of every slide.
    \end{block}
    \begin{alertblock}{Watch out!}
        Never rely on colour alone to carry a meaning.
    \end{alertblock}
    \begin{exampleblock}{Example}
        \texttt{\textbackslash usetheme[accent = 0F766E]\{\ThemeName\}}
    \end{exampleblock}
\end{frame}


\begin{frame}{Tables}
    \scriptsize % I had to go overboard with the table, didn't I? 
    \begin{table}
    \centering
        \begin{tabular}{|l|c|c|c|}
            \toprule
            & Co-efficients & Variables & Constants\\
            \midrule
            Actors & $\displaystyle A={\begin{bmatrix}a_{11}&a_{12}&\cdots &a_{1n}\\a_{21}&a_{22}&\cdots &a_{2n}\\\vdots &\vdots &\ddots &\vdots \\a_{m1}&a_{m2}&\cdots &a_{mn}\end{bmatrix}}$ & $\mathbf {x} ={\begin{bmatrix}x_{1}\\x_{2}\\\vdots \\x_{n}\end{bmatrix}}$ & $\mathbf {b} ={\begin{bmatrix}b_{1}\\b_{2}\\\vdots \\b_{m}\end{bmatrix}}.$\\
            \midrule
            Equation & \multicolumn{3}{c|}{$A\mathbf{x} = \mathbf{b}$}\\
            \midrule
            Story & \multicolumn{3}{c|}{${\displaystyle {\begin{cases}a_{11}x_{1}+a_{12}x_{2}+\dots +a_{1n}x_{n}=b_{1}\\a_{21}x_{1}+a_{22}x_{2}+\do`ts +a_{2n}x_{n}=b_{2}\\\vdots \\a_{m1}x_{1}+a_{m2}x_{2}+\dots +a_{mn}x_{n}=b_{m}\end{cases}}}$}\\
            \bottomrule
        \end{tabular}
        \caption{A table. Good for show, not good as a table.}
    \end{table}
\end{frame}

% Slide demonstrating standout option
\begin{frame}[standout]
    Questions?
\end{frame}

\appendix
\begin{frame}{Backup Slide}
    Appendix slides are not counted for the total.\\
    Package \texttt{appendixnumberbeamer} ensures that.
\end{frame}


% \begin{frame}
    
% \end{frame}


\end{document}